\chapter{뉴런의 유형을 정하지 않는 물질}
\label{sec:othermediators}

\ref{sec:neurontypes}장에서 우리는 뉴런을 주 신경전달물질에 따라 분류해 열 가지 유형을 얻었다(표~\ref{tab:types}). 이제 복잡한 사정을 더할 때이다. 뉴런은 주 신경전달물질만 내보내지 않으며, 뇌에서는 주 신경전달물질 외에 다른 물질도 일한다. 이 물질들은 네트워크가 신호를 전달하고 저장하는 방식을 바꾼다.

주소 버스와 브로드캐스트 버스(\ref{sec:buses}절) 외에 세 번째 버스가 있다. \emph{역방향 채널}이다. 몇몇 물질은 송신자가 아니라 \emph{수신자}가 내보내며, 이 물질은 틈새를 거슬러 송신자의 출력 쪽으로 가서 송신자가 다음번에 신경전달물질을 얼마나 내보낼지를 바꾼다. 통신망으로 치면 송신자가 자기 전송을 조정하는 데 쓰는 수신 확인과 비슷하다. 뉴런이 연결 가중치를 바꿀 때 쓰는 것이 바로 이 채널이다. 이 채널을 통해 송신자의 출력은 수신자의 활동, 즉 \ref{sec:dofa}장 학습 규칙의 두 번째 인자를 알게 된다.

알려진 신경전달물질의 전체 목록은 길다. 100가지가 넘는 물질이 있고, 대부분은 짧은 단백질 사슬인 신경펩타이드이다. 그러나 이들은 모두 몇 개의 부류로 묶이며, 부류 안의 물질은 비슷하게 행동한다. 주 신경전달물질이 되는 일이 없는 물질을 표~\ref{tab:othermediators}에 모았다. 엔지니어가 물을 세 가지 질문, 즉 신호가 어느 버스로 가는지, 빨리 작용하는지, 보통 부호가 무엇인지를 함께 적었다. 이 물질들은 뉴런의 유형을 정하지 않는다. 주 신경전달물질과 함께 방출되거나, 뉴런이 아닌 세포가 방출하거나, 반대 방향으로 가기 때문이다.

\begin{table}[t]
\centering
\footnotesize
\setlength{\tabcolsep}{4pt}
\renewcommand{\arraystretch}{1.12}
\begin{tabular}{>{\raggedright}p{2.25cm}>{\raggedright}p{2.8cm}>{\raggedright}p{1.8cm}>{\raggedright}p{1.5cm}>{\centering}p{0.8cm}>{\raggedright\arraybackslash}p{4.3cm}}
\toprule
부류 & 물질 & 버스 & 속도 & 부호 & 계산에서의 역할 \\
\midrule
아미노산 & D-세린 & 국소 & 느림 & AND & 글루탐산 수용체의 두 번째 열쇠: “AND” 조건 \\
\midrule
모노아민 & 미량 아민 & 브로드캐스트 & 느림 & $\pm$ & 모노아민 채널의 미세 조정 \\
\midrule
\multirow{2}{2.25cm}{퓨린}
 & ATP & 주소 & 빠름과 느림 & $+$ & 공동 신경전달물질. 뉴런과 교세포 사이의 신호 \\
 & 아데노신 & 체적 & 느림 & $-$ & 활동하면 쌓인다: 부하 계수기~\cite{Dunwiddie2001} \\
\midrule
\multirow{8}{2.25cm}{신경펩타이드 (100가지 이상)}
 & 엔케팔린, 엔도르핀, 다이노르핀 & 체적 & 느림 & $-$ & 전달 억제 \\
 & P 물질 & 체적 & 느림 & $+$ & 느린 증폭 \\
 & 신경펩타이드 Y & 체적 & 느림 & $-$ & 느린 억제 \\
 & 소마토스타틴 & 체적 & 느림 & $-$ & 느린 억제, GABA와 함께 나옴 \\
 & 혈관활성 장펩타이드(VIP) & 체적 & 느림 & $+$ & 탈억제: 억제성 뉴런의 억제 \\
 & 콜레시스토키닌 & 체적 & 느림 & $\pm$ & 일부 억제성 뉴런에서 GABA와 함께 나옴 \\
 & 오렉신 & 브로드캐스트 & 느림 & $+$ & 각성 상태의 안정기 \\
 & 옥시토신, 바소프레신 & 브로드캐스트 & 느림 & $\pm$ & 행동의 전역 설정 \\
\midrule
\multirow{2}{2.25cm}{기체}
 & 일산화질소 NO & 역방향, 체적 & 느림 & $\pm$ & 막을 그대로 통과한다. 가중치 변화 때의 역방향 신호~\cite{Garthwaite2008} \\
 & CO, H$_2$S & 체적 & 느림 & $\pm$ & 역할이 거의 연구되지 않음 \\
\midrule
지질 & 엔도카나비노이드 (아난다마이드, 2-AG) & 역방향 & 느림 & $-$ & 수신자가 송신자의 방출을 줄인다: 가중치에 대한 되먹임~\cite{WilsonNicoll2001} \\
\bottomrule
\end{tabular}
\caption{뉴런의 유형을 정하지 않는 물질. \emph{버스}, \emph{속도}, \emph{부호}는 표~\ref{tab:types}와 같다. 여기에 더해, 체적 버스는 신경전달물질이 방출 지점 둘레의 세포 사이 공간으로 퍼지는 것, 역방향은 수신자에서 송신자 쪽으로 거슬러 가는 것, 국소는 방출 지점 가까이에서 작용하는 것이다. “AND”는 허가 신호이다. 그 자체로는 전류를 일으키지 않지만, 이것이 없으면 다른 입력이 작동하지 않는다. 논리 AND 게이트의 두 번째 입력과 같다. 신경펩타이드는 거의 언제나 주 신경전달물질과 함께, 주로 스파이크 빈도가 높을 때 방출된다~\cite{vandenPol2012}.}
\label{tab:othermediators}
\end{table}
